% Part: set-theory
% Chapter: ordinals
% Section: introduction

\documentclass[../../../include/open-logic-section]{subfiles}

\begin{document}

\olfileid{sth}{ordinals}{intro}

\olsection{جان پہچان}

\olref[z][]{chap} وچ اسیں \stagesinf{} دے روپ وچ فرض کیتا سی کہ
درجہ بندی دا اک لامتناہی مرحلہ اے (ساڈا لامتناہیت دا مسلّمہ وی ویکھو)۔
پر \stagessucc{} مطابق اسیں لامتناہی مرحلے اُتے رُک نہیں سکدے۔
سانوں اگے ودھنا پئے گا۔ سو پہلے لامتناہی مرحلے توں اگلے مرحلے اُتے
اسیں اوہناں سیٹاں دے سارے ممکن اکٹھ بناندے آں جیہڑے پہلے لامتناہی
مرحلے اُتے موجود سن۔ تے فیر ایہو کم کردے آں، تے فیر، تے فیر، \dots

ایس وچ ان کہی گل ایہ اے کہ اسیں عدد دا اک ’’سوجھ والا‘‘ تصور ورتنا
شروع کر دتا اے، جیہدے مطابق سارے قدرتی عدداں توں \emph{پچھوں} وی عدد
ہو سکدے نیں۔ خاص کر کے ایہ \emph{لامتناہی درجے دے ترتیبی عدد} دا خیال
اے۔ ایس باب دا مقصد ایس خیال نوں ہور ٹھیک بنانا اے۔ اسیں ترتیبی عدد
دا عام تصور ویکھاں گے تے فیر کجھ خاص سیٹاں نوں اپنے ترتیبی عدد دے طور
اُتے صاف تعریف کراں گے۔

\end{document}
